\section{Single-calibration closure and cross-observable predictions}
\label{sec:bhsm-single-calibration}

The preceding theorems reduce the cosmological anisotropic response to one
physical $n=2$ phase-space state $X_2$ and one common response operator
$\mathcal R_2$.  The remaining overall normalization of the linear one-mode
branch can therefore be fixed by one observational calibration without
introducing an independent amplitude for each observable.

\subsection{Amplitude factorization}

For a fixed background branch and a fixed normalized physical mode shape,
write
\begin{equation}
X_2(z)=A_2\,\widehat X_2(z),
\end{equation}
where $\widehat X_2$ is the unit-normalized reduced solution and $A_2$ is the
single branch amplitude.  For any scalar linear observable,
Eq.~\eqref{eq:final-observable-kernel} becomes
\begin{equation}
\mathcal A_{\mathcal O}(z)
=
A_2\,\mathcal F_{\mathcal O}(z),
\label{eq:observable-F-kernel}
\end{equation}
with
\begin{equation}
\mathcal F_{\mathcal O}(z)
=
\mathcal L_{\mathcal O,z}
\mathcal R_2(z)
\widehat X_2(z).
\end{equation}
The functions $\mathcal F_{\mathcal O}$ are fixed by the reduced action and
the chosen background branch.

\subsection{Single-calibration theorem}

\paragraph{Theorem 4 (one calibration fixes all linear amplitudes).}
Suppose there exists one calibration observable $\mathcal C$ at redshift
$z_\star$ for which
$\mathcal F_{\mathcal C}(z_\star)\neq0$.  Then the single-mode amplitude is
\begin{equation}
\boxed{
A_2
=
\frac{
\mathcal A_{\mathcal C}^{\rm obs}(z_\star)
}{
\mathcal F_{\mathcal C}(z_\star)
}
}
\label{eq:A2-calibration}
\end{equation}
and every other linear observable is predicted by
\begin{equation}
\boxed{
\mathcal A_{\mathcal O}(z)
=
\mathcal A_{\mathcal C}^{\rm obs}(z_\star)
\frac{
\mathcal F_{\mathcal O}(z)
}{
\mathcal F_{\mathcal C}(z_\star)
}.
}
\label{eq:cross-observable-prediction}
\end{equation}
Thus no second anisotropic amplitude parameter is available once the branch,
mode normalization and one calibration are fixed.

\paragraph{Proof.}
Equation~\eqref{eq:observable-F-kernel} is linear in the single branch
amplitude $A_2$.  Solving the calibration equation for $A_2$ gives
Eq.~\eqref{eq:A2-calibration}.  Substitution into the response of any other
observable gives Eq.~\eqref{eq:cross-observable-prediction}. \hfill$\square$

\subsection{Reference supernova calibration}

The current reference branch carries forward the low-redshift supernova
distance-modulus calibration
\begin{equation}
A_\mu^\star
=
-0.0412\pm0.009,
\qquad
0.01<z<0.03,
\label{eq:reference-A-mu}
\end{equation}
with reference sky axis
\begin{equation}
{\rm RA}=211.48^\circ,
\qquad
{\rm Dec}=-12.81^\circ .
\label{eq:reference-axis}
\end{equation}
Because
\begin{equation}
\delta\mu
=
\frac{5}{\ln 10}
\frac{\delta D_L}{D_L},
\end{equation}
this corresponds to
\begin{equation}
\boxed{
\left(
\frac{\delta D_L}{D_L}
\right)_\star
=
-0.01897\pm0.00414
}
\label{eq:reference-DL-fraction}
\end{equation}
for the reference calibration bin.

Equations~\eqref{eq:reference-A-mu}--\eqref{eq:reference-DL-fraction} are
calibration inputs to the reference branch.  They are not derived by
Theorem~4 and they are not treated as independent calibrations of the other
observables.

\subsection{Parameter-free amplitude ratios after calibration}

For any two observables on the same frozen branch,
\begin{equation}
\boxed{
\frac{
\mathcal A_{\mathcal O_1}(z_1)
}{
\mathcal A_{\mathcal O_2}(z_2)
}
=
\frac{
\mathcal F_{\mathcal O_1}(z_1)
}{
\mathcal F_{\mathcal O_2}(z_2)
}
}
\label{eq:amplitude-ratio}
\end{equation}
provided the denominator is nonzero.  The overall mode amplitude cancels.

This is the principal route from a one-observable calibration to genuine
cross-observable prediction.  In particular, once the numerical kernels are
evaluated, the supernova calibration fixes the absolute anisotropic
normalization for
\begin{equation}
\delta H/H,\qquad
\delta D_A/D_A,\qquad
\delta(f\sigma_8)/(f\sigma_8),\qquad
\delta(\Phi+\Psi),
\end{equation}
without fitting those amplitudes independently.

\subsection{Shared-axis theorem as a falsification test}

The pure one-mode branch also carries one common axis
$\hat{\mathbf p}$.  Therefore
\begin{equation}
\delta\mathcal O_i(z,\hat{\mathbf n})
=
\mathcal A_{\mathcal O_i}(z)
(\hat{\mathbf n}\!\cdot\!\hat{\mathbf p})
\end{equation}
for every scalar linear observable in the branch.

A persistent, statistically resolved disagreement between the axes inferred
from independent observables, after survey masks, calibration systematics,
and local-structure contamination are controlled, falsifies the pure
single-$n=2$ realization even if individual amplitudes agree.

\subsection{Reference-branch versus framework-wide statements}

Theorem~4 is structural.  The numerical functions
$\mathcal F_{\mathcal O}(z)$ are branch-dependent.  Accordingly:

\begin{itemize}
\item the existence of one common amplitude $A_2$ and one common axis is a
      consequence of the single-mode linear branch;
\item the particular numerical values of
      $\mathcal F_{\mathcal O}(z)$ on R1 are reference-branch predictions;
\item the supernova amplitude in
      Eq.~\eqref{eq:reference-A-mu} is an observational calibration input;
\item no cross-observable amplitude may be promoted to a framework-wide
      prediction until its kernel has been evaluated from the reduced
      operator chain.
\end{itemize}

\subsection{Closure target}

The final numerical closure task for the manuscript is now finite and
explicit:
\begin{equation}
\boxed{
\left\{
\mathcal F_{D_L}(z),
\mathcal F_H(z),
\mathcal F_{D_A}(z),
\mathcal F_{f\sigma_8}(z),
\mathcal F_{\rm lens}(z)
\right\}_{\rm R1}.
}
\label{eq:final-kernel-set}
\end{equation}
Once the five frozen-branch kernels in
Eq.~\eqref{eq:final-kernel-set} are evaluated from the already-defined
operators, one supernova calibration determines all corresponding
anisotropic amplitudes through
Eq.~\eqref{eq:cross-observable-prediction}.
